% Argumenten las principales diferencias observadas en las dinámicas del pozo nito simétrico y no simétrico. Les recomiendo que varíen los parámetros y analicen diferentes situaciones que puedan ocurrir.

% ¿Cuáles conceptos aprendimos con este proyecto? ¿Cuáles habilidades desarrollamos con el desarrollo del proyecto?

\subsection{Andros Izhar Alvarado Sánchez}

Al comparar los dos pozos, me llamó la atención cómo la simetría se refleja en sus estados y en su evolución. En el simétrico, los estados tienen paridad definida; en el asimétrico, las paredes diferentes hacen que la función de onda se comporte distinto a cada lado. Con los parámetros del proyecto encontramos 11 estados atrapados en el primer caso y 9 en el segundo.

Lo que más me sirvió fue ver que las gráficas tenían sentido a partir de lo que habíamos planteado en el modelo. Al programarlo hubo que revisar varias veces que los resultados fueran razonables; eso me hizo fijarme más en cómo pequeños cambios en el potencial pueden cambiar lo que vemos en la simulación.
